% Generated draft: requires mathematical and source review.
\documentclass[11pt]{article}
\usepackage[margin=25mm]{geometry}
\usepackage{fontspec}
\setmainfont{Latin Modern Roman}
\usepackage{amsmath,amssymb,mathtools,unicode-math}
\setmathfont{Latin Modern Math}
\usepackage{enumitem,needspace,xurl,hyperref}
\setlength{\parindent}{0pt}
\setlength{\parskip}{5pt}
\setlength{\emergencystretch}{3em}
\newcommand{\N}{\mathbb{N}}
\newcommand{\Z}{\mathbb{Z}}
\newcommand{\Q}{\mathbb{Q}}
\newcommand{\R}{\mathbb{R}}
\newcommand{\C}{\mathbb{C}}
\newcommand{\F}{\mathbb{F}}
\newcommand{\A}{\mathbb{A}}
\newcommand{\PP}{\mathbb{P}}
\newcommand{\E}{\mathbb{E}}
\newcommand{\eps}{\varepsilon}
\newcommand{\dd}{\,\mathrm{d}}
\begin{document}
\section*{UnsolvedMath ID: 30005576}
\textbf{Source problem code:} catalogue entry
\subsection{Definitions and mathematical statement}
\textbf{Unrecoverable mathematical statement.} The supplied record for permanent UnsolvedMath ID \(30005576\), source code \texttt{catalogue entry}, is titled ``Geometric Analysis Table-of-Contents Fragment.'' Its statement consists of table-of-contents entries, beginning with the isolated word ``Conjecture'' and page number 1627, followed by talk titles, contributors, and page numbers. It supplies no conjecture text, mathematical definitions, hypotheses, domains, boundary conditions, quantifiers, or conclusion from which a precise mathematical proposition can be recovered. The adjacent talk titles do not determine the intended problem or the attribution of the isolated conjecture entry.

Consequently, no mathematical assertion is formulated here. In particular, neither a conjecture about shrinking gradient Kähler--Ricci solitons nor a problem associated with any other listed talk can be inferred from this fragment. Recovering an intended target requires the relevant complete problem passage and its attribution.

The supplied registry marks the statement as \texttt{unrecoverable}, provides no clean statement, and leaves the proposer unspecified. Its verification field reports that the official report identifies the extract as a table-of-contents fragment rather than a mathematical question, and that the original record remains publicly visible for provenance. This is supplied verification information, not an independent verification performed for this draft. The registry's \texttt{open} label is expressly a conservative curation classification because the fragment does not state a determinate conjecture; it does not establish the existence or unresolved status of a particular mathematical proposition.
\subsection{Short English statement}
Permanent UnsolvedMath ID \(30005576\) preserves a table-of-contents fragment, not a recoverable mathematical question. The isolated ``Conjecture'' entry has no statement or identifiable proposer. A precise formulation requires the missing source passage. The supplied \texttt{open} classification is dated curation rather than a verified mathematical status.
\subsection{Sources}
\begin{itemize}
\item Supplied local UnsolvedMath registry record: permanent ID \(30005576\), exact source code \texttt{catalogue entry}, title ``Geometric Analysis Table-of-Contents Fragment,'' collection ``Oberwolfach Reports Open Problems.'' The collection is described as containing open and partially solved problems extracted from workshop discussions, with report provenance and dated literature assessments. The supplied background describes Oberwolfach Reports as documenting mathematical workshops and their open-problem discussions.
\item Supplied report citation: \emph{Differentialgeometrie im Grossen (2024)}, Oberwolfach Reports, DOI \texttt{10.4171/owr/2023/29}, pp.\ 1617--1670. Supplied source URL and literature-source link, labeled ``original OWR report'': \url{https://doi.org/10.4171/owr/2023/29}. This is the cited report; no complete authoritative problem passage is supplied. The record gives proposed year 2023 and the tag \texttt{owr-year:2023}, whereas the citation labels the report 2024; this discrepancy is unresolved.
\item Supplied literature assessment, checked on \textbf{2026-08-22}: the extract is a table-of-contents fragment without a determinate mathematical conjecture, so no resolution can be established from the record and it is conservatively classified as open. This is supplied dated curation, not independently verified. The supplied background requests independent verification of time-sensitive status and cited claims. The supplied statement-verification field likewise reports confirmation from the official report; this draft does not claim to have visited that report.
\item Fragment entry: ``Conjecture,'' p.\ 1627. No associated statement or identifiable proposer is supplied.
\item Fragment attribution: Ronan J. Conlon, joint with Richard Bamler, Charles Cifarelli, and Alix Deruelle, ``Two-dimensional shrinking gradient Kähler-Ricci solitons,'' p.\ 1628.
\item Fragment attribution: Hokuto Konno, joint with Jianfeng Lin, ``Seiberg-Witten theory and moduli spaces of 4-manifolds and metrics,'' p.\ 1631.
\item Fragment attribution: Christos Mantoulidis, joint with Otis Chodosh and Felix Schulze, ``Improved generic regularity of minimizing hypersurfaces,'' p.\ 1634.
\item Fragment attribution: Bin Guo, joint with Duong H. Phong, Jian Song, and Jacob Sturm, ``Diameter estimates in Kähler geometry,'' p.\ 1636.
\item Fragment attribution: Claude LeBrun, ``Einstein Manifolds, Weyl Curvature, \& Conformally Kähler Geometry,'' p.\ 1638.
\item Fragment attribution: Carla Cederbaum, joint with Melanie Graf and Jan Metzger, ``Coordinates are messy in General (Relativity),'' p.\ 1640.
\item Fragment attribution: Paula Burkhardt-Guim, ``ADM mass for \(C^0\) metrics and distortion under Ricci-DeTurck flow,'' p.\ 1643. The extracted title renders the regularity notation as ``C 0.''
\item Fragment attribution: Daniel Stern, joint with Davide Parise and Alessandro Pigati, ``Minimal submanifolds and mean curvature flows of codimension 2 via adiabatic limits,'' p.\ 1646.
\item Fragment attribution: Cristiano Spotti, joint with Daniele Angella, Simone Calamai, and Francesco Pediconi, ``Towards canonical locally conformally Kähler metrics,'' p.\ 1649.
\item Fragment attribution: Robin Neumayer, ``Epsilon regularity for spaces with scalar curvature lower bounds,'' p.\ 1649.
\item Fragment attribution: Jonathan Rosenberg, joint with Shmuel Weinberger, ``Positive scalar curvature on manifolds with boundary and their doubles,'' p.\ 1651.
\item Fragment attribution: Roman Sauer, joint with Sabine Braun, ``Volume and macroscopic scalar curvature,'' p.\ 1653. The fragment ends with ``Workshop: Differentialgeometrie im Grossen Table of Contents.'' These contributor attributions belong to the listed entries and do not identify a proposer of the isolated conjecture.
\end{itemize}
\end{document}
